\documentclass[11pt,a4paper]{article}
\usepackage[utf8]{inputenc}
\usepackage[italian]{babel}
\usepackage{geometry}
\geometry{top=2.5cm,bottom=2.5cm,left=2.5cm,right=2.5cm}
\usepackage{amsmath,amssymb}
\usepackage{tcolorbox}
\tcbuselibrary{skins,breakable}
\usepackage{booktabs}
\usepackage{tabularx}
\usepackage{enumitem}
\usepackage{xcolor}
\usepackage{listings}
\usepackage{hyperref}
\hypersetup{
    colorlinks=true,
    linkcolor=blue!70!black,
    citecolor=blue!70!black,
    urlcolor=teal!70!black
}

\lstset{
    basicstyle=\ttfamily\small,
    breaklines=true,
    frame=single,
    backgroundcolor=\color{gray!10},
    keywordstyle=\color{blue}\bfseries,
    commentstyle=\color{green!50!black}\itshape,
    stringstyle=\color{red!70!black}
}

\newtcolorbox{questionbox}[1][]{
    colback=gray!5!white,
    colframe=gray!60!black,
    fonttitle=\bfseries,
    title={#1},
    breakable
}

\newtcolorbox{solutionbox}[1][]{
    colback=blue!5!white,
    colframe=blue!70!black,
    fonttitle=\bfseries,
    title={#1},
    breakable
}

\newtcolorbox{altbox}[1][]{
    colback=orange!5!white,
    colframe=orange!80!black,
    fonttitle=\bfseries,
    title={#1},
    breakable
}

\begin{document}

\title{\textbf{Sicurezza dei Sistemi Informativi} \\ \large Appello d'Esame del 27 Gennaio 2023 \\ \small Anno Accademico 2022-2023}
\author{Politecnico di Torino --- Docente: Prof. Cataldo Basile}
\date{2023-01-27}
\maketitle


\section{Testo Ufficiale dell'Esame}

\begin{questionbox}[Domanda 1: Analisi Ciphersuite TLS 1.2]
Un client e un server del protocollo TLS 1.2 negoziano la seguente ciphersuite:
\begin{center}
\texttt{TLS\_DHE\_RSA\_WITH\_3DES\_EDE\_CBC\_SHA256}
\end{center}
\begin{enumerate}[label=\alph*.]
    \item Spiegare cosa significa ciascuna delle componenti principali della ciphersuite e su quale propriet\`a di sicurezza (riservatezza, integrit\`a, autenticazione, segretezza in avanti) ha influenza.
    \item Spiegare brevemente cosa accade durante il protocollo di Handshake per lo scambio delle chiavi e l'autenticazione del server.
    \item Indicare quale \`e la parte pi\`u critica e obsoleta di questa ciphersuite e per quali motivi tecnici e vulnerabilit\`a note (es. Sweet32, problemi di CBC).
\end{enumerate}
\end{questionbox}

\begin{questionbox}[Domanda 2: Verifica di Firma Digitale e Certificato X.509]
Bob e Alice sono colleghi presso la societ\`a SSI Inc. Alice invia a Bob:
\begin{itemize}
    \item Il suo certificato X.509 a chiave pubblica firmato dalla Certification Authority (CA) aziendale;
    \item Un file di dati firmato digitalmente con la chiave privata corrispondente alla chiave pubblica presente nel certificato.
\end{itemize}
Assumendo che la SSI Inc sia dotata di una Public Key Infrastructure (PKI) completa e gestita correttamente sia a livello di registrazione che di infrastruttura, descrivere dettagliatamente tutti i passi che Bob (o un software di verifica al suo posto) deve compiere per stabilire con certezza che la firma sia pienamente valida.
\end{questionbox}

\begin{questionbox}[Domanda 3: Privacy by Design e il Ruolo del DPO nel GDPR]
Riportare e spiegare 4 aspetti o principi operativi relativi all'implementazione di servizi conformi al principio di \textbf{Privacy by Design} secondo il GDPR.
Inoltre, spiegare chi \`e il \textbf{Data Protection Officer (DPO)}, quali sono i suoi compiti istituzionali e che tipo di responsabilit\`a giuridica possiede in caso di data breach.
\end{questionbox}

\begin{questionbox}[Domanda 4: Autenticazione basata su Password Ripetibili]
Spiegare il funzionamento dell'autenticazione basata su password ripetibili (memorizzate e riutilizzabili).
Descrivere le principali problematiche di sicurezza che bisogna affrontare:
\begin{itemize}
    \item Lato client;
    \item Nel canale di comunicazione;
    \item Lato server (memorizzazione);
    \item Problematiche generali e comportamentali degli utenti.
\end{itemize}
Per ciascuna problematica, indicare le relative misure tecniche di mitigazione (KDF, salting, canali protetti, MFA).
\end{questionbox}

\begin{questionbox}[Domanda 5: Meccanismi Anti-SYN Flooding: SYN Interceptor vs SYN Monitor]
Indicare quali delle seguenti affermazioni su \textbf{SYN Interceptor} e \textbf{SYN Monitor} sono vere:
\begin{enumerate}[label=\alph*.]
    \item Il SYN Interceptor \`e meno efficace del SYN Monitor;
    \item Il SYN Monitor \`e meno efficace del SYN Interceptor;
    \item Non alterano il three-way handshake ma inviano pacchetti al posto del server;
    \item Cambiano il three-way handshake per evitare SYN flooding inviando l'hash del sequence number;
    \item Il SYN Interceptor risponde al posto del server. Contatta il server stabilendo una connessione solo se il three-way handshake ha successo;
    \item Servono entrambi a mitigare gli attacchi di tipo SYN flooding;
    \item Il SYN Monitor risponde al posto del server. Contatta il server stabilendo una connessione solo se il three-way handshake ha successo;
    \item Nessuno dei due risolve completamente il problema del SYN flooding.
\end{enumerate}
\end{questionbox}

\begin{questionbox}[Domanda 6: Progettazione Architettura di Rete e Firewalling]
Un'azienda ha una rete suddivisa in tre sottoreti interne:
\begin{itemize}
    \item Sottorete 1: \texttt{192.168.1.0/24} (ospita server interno \texttt{192.168.1.1});
    \item Sottorete 2: \texttt{192.168.2.0/24} (ospita server interno \texttt{192.168.2.3});
    \item Sottorete 3: \texttt{192.168.3.0/24};
    \item Server pubblici con indirizzi IP pubblici esposti su Internet.
\end{itemize}
Politica di sicurezza richiesta:
\begin{enumerate}
    \item Nessuna restrizione interna ad ognuna delle 3 sottoreti;
    \item Il server \texttt{192.168.1.1} deve essere raggiungibile dalle sottoreti .2 e .3;
    \item Il server \texttt{192.168.2.3} deve essere raggiungibile dalla sottorete .1;
    \item Nessuna sottorete interna (.1, .2, .3) deve poter essere raggiunta dall'esterno (Internet);
    \item Gli host delle sottoreti .1 e .2 possono raggiungere Internet;
    \item Gli host della sottorete .3 non possono raggiungere Internet;
    \item I server pubblici devono essere raggiungibili da Internet, ma solo su un insieme di URL note (es. \texttt{\$url1, \$url2, ...}).
\end{enumerate}
Indicare quale architettura e quali controlli preventivi abilitare sull'elemento al bordo e riportare le corrispondenti regole di configurazione.
\end{questionbox}

\newpage
\section{Soluzione Dettagliata e Guida Risolutiva}

\begin{solutionbox}[Soluzione Domanda 1: Ciphersuite TLS 1.2]
\textbf{a. Decodifica delle Componenti della Ciphersuite:}
\begin{itemize}
    \item \texttt{TLS}: Specifica il protocollo di appartenenza.
    \item \texttt{DHE} (Diffie-Hellman Ephemeral): Meccanismo di scambio delle chiavi. Genera un segreto condiviso effimero garantendo la propriet\`a di \textbf{Perfect Forward Secrecy (PFS)}.
    \item \texttt{RSA}: Algoritmo di crittografia asimmetrica utilizzato dal server esclusivamente per \textbf{autenticarsi} presso il client (firmando digitalmente i parametri DH effimeri).
    \item \texttt{3DES\_EDE\_CBC}: Algoritmo e modalit\`a di cifratura simmetrica per garantire la \textbf{confidenzialit\`a} dei dati. Triplo DES con tre fasi (Encrypt-Decrypt-Encrypt) in modalit\`a Cipher Block Chaining.
    \item \texttt{SHA256}: Funzione di hash crittografica a 256 bit utilizzata sia come generatore di MAC (HMAC-SHA256) per garantire \textbf{integrit\`a e autenticit\`a} dei pacchetti, sia come base per la Pseudo-Random Function (PRF) per derivare le chiavi.
\end{itemize}

\textbf{b. Fasi dell'Handshake per Scambio Chiavi e Autenticazione:}
Il server genera parametri Diffie-Hellman effimeri casuali $(g, p, g^b \pmod p)$ e li firma digitalmente con la propria chiave privata RSA. Nel messaggio \texttt{ServerKeyExchange}, invia tali parametri assieme alla firma e al proprio certificato X.509 (\texttt{Certificate}). Il client valida il certificato risalendo alla CA, verifica la firma RSA sui parametri DH usando la chiave pubblica del certificato (autenticazione del server avvenuta). Successivamente, il client genera il proprio valore effimero $g^a \pmod p$, lo invia nel messaggio \texttt{ClientKeyExchange}, e calcola il Pre-Master Secret $g^{ab} \pmod p$.

\textbf{c. Criticit\`a della Ciphersuite:}
La componente pi\`u critica \`e senza dubbio \textbf{\texttt{3DES\_EDE\_CBC}}:
\begin{enumerate}
    \item \textbf{Vulnerabilit\`a Sweet32 (CVE-2016-2183)}: 3DES utilizza una dimensione di blocco obsoleta di soli \textbf{64 bit} (8 byte). A causa del paradosso del compleanno, dopo che una connessione TLS ha cifrato circa $2^{32}/2 = 2^{32}$ blocchi ($\approx 32$ GB di dati) sotto la stessa chiave, si verifica una collisione nei blocchi di ciphertext con probabilit\`a superiore al 50\%. In modalit\`a CBC, ci\`o consente a un attaccante Man-in-the-Middle (sfruttando codice JavaScript malevolo che induce il browser a generare richieste ripetute) di estrarre cookie HTTP di sessione in chiaro.
    \item \textbf{Padding Oracle Attack (Vaudenay)}: La combinazione MAC-then-Encrypt con modalit\`a CBC \`e storicamente affetta da attacchi a canale laterale sui messaggi di errore di padding (Lucky Thirteen, POODLE).
    \item \textbf{Prestazioni e Chiave Efficace}: 3DES \`e estremamente lento in software e la sua sicurezza effettiva \`e di soli 112 bit (invece di 168 bit) a causa degli attacchi meet-in-the-middle. \`E oggi ufficialmente deprecato dal NIST.
\end{enumerate}
\end{solutionbox}

\begin{solutionbox}[Soluzione Domanda 2: Catena di Verifica del Certificato e della Firma Digitale]
Per accertare la validit\`a legale e tecnica del documento e del certificato, Bob deve compiere i seguenti passi sequenziali e rigorosi:
\begin{enumerate}
    \item \textbf{Verifica del Periodo di Validit\`a}: Bob verifica che il timestamp corrente rientri nell'intervallo temporale specificato nei campi \texttt{notBefore} e \texttt{notAfter} del certificato di Alice.
    \item \textbf{Costruzione e Verifica della Catena di Fiducia (Trust Chain)}: Bob estrae il nome dell'emittente (\texttt{Issuer}) dal certificato di Alice e recupera il certificato della CA aziendale SSI Inc. Verifica che la firma digitale apposta sul certificato di Alice sia valida ricalcolando il digest del certificato e verificandolo con la chiave pubblica della CA. Se la CA aziendale \`e intermedia, il processo si ripete ricorsivamente fino a raggiungere un certificato radice autofirmato (\textit{Root CA}) presente nel Trust Store locale fidato di Bob.
    \item \textbf{Verifica dello Stato di Revoca}: Bob deve accertarsi che il certificato di Alice non sia stato revocato prima della sua scadenza naturale (es. per compromissione della chiave privata o cessazione del rapporto lavorativo). Ci\`o avviene tramite:
    \begin{itemize}
        \item Download e verifica della \textbf{Certificate Revocation List (CRL)} firmata dalla CA;
        \item Oppure tramite interrogazione in tempo reale via protocollo \textbf{OCSP} (Online Certificate Status Protocol) o verifica del timestamp OCSP Stapling.
    \end{itemize}
    \item \textbf{Verifica delle Estensioni e del Key Usage}: Bob controlla che l'estensione X.509 \texttt{KeyUsage} e/o \texttt{ExtendedKeyUsage} autorizzi esplicitamente l'uso della chiave per la firma digitale (\textit{Digital Signature}, \textit{Non-Repudiation}).
    \item \textbf{Verifica della Firma Digitale sul Documento}:
    \begin{enumerate}
        \item Bob calcola autonomamente il digest crittografico $H(M)$ del file ricevuto utilizzando la stessa funzione di hash specificata nell'algoritmo di firma (es. SHA-256).
        \item Decifra la firma digitale $S$ utilizzando la \textbf{chiave pubblica di Alice} estratta dal certificato validato (ad esempio applicando la verifica RSA $S^e \pmod n$ o la formula di verifica ECDSA).
        \item Confronta il digest ottenuto dalla firma con il digest calcolato $H(M)$. Se coincidono perfettamente, \`e matematicamente garantita l'\textbf{integrit\`a} del file (nessuna modifica post-firma) e l'\textbf{autenticit\`a/non ripudio} di Alice.
    \end{enumerate}
\end{enumerate}
\end{solutionbox}

\begin{solutionbox}[Soluzione Domanda 3: Privacy by Design e DPO]
\textbf{Principio di Privacy by Design (Art. 25 GDPR):}
Prevede che la protezione dei dati personali sia integrata fin dalla fase di progettazione e concezione architetturale di qualsiasi sistema, servizio o processo aziendale. I 4 aspetti implementativi fondamentali includono:
\begin{enumerate}
    \item \textbf{Minimizzazione e Pseudonimizzazione}: Il servizio deve trattare e collezionare solo i dati strettamente necessari. Identificativi diretti (es. codice fiscale) devono essere sostituiti con identificativi pseudonimi o token casuali.
    \item \textbf{Cifratura pervasiva (At Rest e In Transit)}: I dati memorizzati su database/dischi devono essere cifrati (es. AES-256) e le comunicazioni devono essere protette tramite TLS moderno con mutua autenticazione.
    \item \textbf{Privacy by Default}: Le impostazioni predefinite di sicurezza e riservatezza del servizio devono essere configurate al massimo livello di protezione, senza richiedere azioni correttive all'utente.
    \item \textbf{Gestione del Ciclo di Vita e Conservazione}: Implementazione di meccanismi automatici di scadenza e cancellazione sicura (\textit{data retention}) al termine del periodo necessario.
\end{enumerate}

\textbf{Ruolo e Responsabilit\`a del Data Protection Officer (DPO):}
Il DPO \`e una figura di garanzia indipendente. I suoi compiti principali sono:
\begin{itemize}
    \item Consulenza e informazione al Titolare e ai dipendenti sugli obblighi del GDPR;
    \item Sorveglianza sull'osservanza del Regolamento e delle policy aziendali di sicurezza;
    \item Cooperazione e punto di contatto con l'Autorit\`a Garante per la Privacy.
\end{itemize}
\textit{Responsabilit\`a legale}: Il DPO **non \`e legalmente responsabile** per eventuali violazioni o data breach dell'azienda e non risponde dei danni agli interessati. La responsabilit\`a civile, penale e amministrativa (sanzioni fino a 20 milioni di euro o al 4\% del fatturato annuo globale) ricade esclusivamente sul **Titolare del Trattamento**.
\end{solutionbox}

\begin{solutionbox}[Soluzione Domanda 4: Autenticazione basata su Password Ripetibili]
\textbf{Funzionamento:} L'utente presenta un identificativo univoco (username) e un segreto condiviso a priori (password statica memorizzata) per dimostrare la propria identit\`a.

\textbf{Problematiche e Mitigazioni per Ambito:}
\begin{center}
\small
\begin{tabularx}{\textwidth}{lXX}
\toprule
\textbf{Ambito} & \textbf{Vulnerabilit\`a / Minacce} & \textbf{Mitigazioni Tecniche} \\
\midrule
\textbf{Lato Client} & Keylogger, malware, spyware, shoulder surfing, memorizzazione non sicura nei browser. & Virtual keyboard, isolamento della memoria, password manager con crittografia forte, anti-malware ed EDR. \\
\textbf{Canale di Rete} & Eavesdropping, packet sniffing, attacchi Man-in-the-Middle e Replay Attack. & Cifratura di canale obbligatoria con TLS 1.3 con certificati verificati (HTTPS), HSTS, protocolli Challenge-Response. \\
\textbf{Lato Server} & Data breach del database delle credenziali; attacchi a dizionario offline e tabelle precalcolate (Rainbow Tables). & \textbf{Mai memorizzare password in chiaro}. Uso di Key Derivation Functions (KDF) memory-hard con sale crittografico (\textit{Salt}) univoco (es. Argon2id, bcrypt, scrypt, PBKDF2) e numero elevato di iterazioni. \\
\textbf{Comportamento} & Scelta di password banali o corte, riutilizzo della stessa password su pi\`u siti (\textit{Credential Stuffing}). & Policy di complessit\`a e lunghezza (minimo 12-16 caratteri), verifica contro elenchi di password compromesse (es. HaveIBeenPwned), \textbf{Multi-Factor Authentication (MFA)} obbligatoria (FIDO2/WebAuthn, TOTP). \\
\bottomrule
\end{tabularx}
\end{center}
\end{solutionbox}

\begin{solutionbox}[Soluzione Domanda 5: SYN Interceptor vs SYN Monitor]
Analisi puntuale ed esaustiva di ciascuna affermazione:
\begin{itemize}
    \item \textbf{a. FALSA}: Il SYN Interceptor \`e \textbf{pi\`u efficace} (e non meno efficace) del SYN Monitor poich\'e intercetta e termina preventivamente il three-way handshake al posto del server, proteggendo la coda delle connessioni semi-aperte (\textit{backlog queue}) del server protetto.
    \item \textbf{b. VERA}: Il SYN Monitor \`e \textbf{meno efficace} del SYN Interceptor perch\'e opera passivamente: non intercetta i pacchetti SYN n\'e risponde al posto del server; osserva solo il traffico e, se nota che una connessione semi-aperta non riceve l'ACK finale entro una soglia temporale, invia un pacchetto \texttt{RST} forzato al server per liberare lo slot. Nel frattempo, il server subisce l'occupazione della memoria.
    \item \textbf{c. FALSA}: Entrambi i meccanismi alterano o intervengono attivamente nella gestione del three-way handshake (il SYN Interceptor si sostituisce direttamente al server inviando SYN-ACK al client; il SYN Monitor inietta pacchetti RST postumi).
    \item \textbf{d. FALSA}: La tecnica che modifica la sequenza numerica inviando un hash crittografico nel sequence number iniziale (ISN) \`e quella dei \textbf{SYN Cookies}, un meccanismo interno allo stack TCP dell'OS che non c'entra con apparati di monitoraggio/intercettazione perimetrali.
    \item \textbf{e. VERA}: Il SYN Interceptor opera a tutti gli effetti come un proxy TCP: risponde al SYN del client con un SYN-ACK forgiato; solo dopo che il client ha risposto con l'ACK valido (comprovando che l'indirizzo IP sorgente non \`e spoofato ed \`e legittimo), contatta il server backend stabilendo una connessione reale.
    \item \textbf{f. VERA}: Entrambi i sistemi sono stati esplicitamente ideati e standardizzati come contromisure perimetrali contro gli attacchi DoS di tipo TCP SYN Flooding.
    \item \textbf{g. FALSA}: L'affermazione descrive il funzionamento del SYN Interceptor, non del SYN Monitor (il quale non risponde mai al posto del server).
    \item \textbf{h. VERA}: Nessuno dei due risolve completamente e definitivamente il problema: il SYN Monitor lascia il server esposto a sovraccarico momentaneo, mentre il SYN Interceptor sposta su di s\'e l'intero onere dello stato TCP, diventando esso stesso potenziale bersaglio di DoS/esaurimento risorse di memoria o collo di bottiglia prestazionale.
\end{itemize}
\end{solutionbox}

\begin{solutionbox}[Soluzione Domanda 6: Architettura di Rete e Tabella Regole Firewall]
\textbf{Architettura Selezionata e Topologia Screened Subnet:}
Per soddisfare tutti i requisiti, in particolare il filtraggio selettivo sui server pubblici basato su un insieme di URL note (\texttt{\$url1, \$url2, ...}), un packet filter di livello 3/4 non \`e sufficiente: il router/firewall di bordo deve instradare il traffico HTTP/HTTPS a un **Reverse Proxy L7 (WAF)** posizionato in DMZ.
\begin{center}
\begin{verbatim}
                           +--------------------+
                           |      INTERNET      |
                           +---------+----------+
                                     | eth0 (IP Pubblico)
                        +------------v------------+
                        |    FIREWALL STATEFUL    |
                        +---+--------+--------+---+
          eth1 (.1)         |        |        |        eth3 (.3)
       192.168.1.1/24       |        |        |      192.168.3.1/24
    +-----------------------+        |        +-----------------------+
    |                                |                                |
+---v------------------+  +----------v---------+  +-------------------v--+
|     SOTTORETE 1      |  |    DMZ / SVR PUB   |  |     SOTTORETE 3      |
|    192.168.1.0/24    |  |   192.168.100.0/24 |  |    192.168.3.0/24    |
|                      |  |                    |  |                      |
| [Server 1] .1.1      |  | [Reverse Proxy].10 |  | [Client Isolati]     |
| [Host Interni]       |  | [Pub Server]   .20 |  | (No Internet)        |
+----------------------+  +--------------------+  +----------------------+
                                eth2 (.2)
                              192.168.2.1/24
                       +-------------v-------------+
                       |        SOTTORETE 2        |
                       |       192.168.2.0/24      |
                       |                           |
                       |   [Server 2] .2.3         |
                       |   [Host Interni]          |
                       +---------------------------+
\end{verbatim}
\end{center}

\textbf{Subnetting, Instradamento e Logica di Filtraggio:}
\begin{itemize}
    \item \textbf{Instradamento Internet}: Le sottoreti .1 e .2 escono su Internet tramite Source NAT (SNAT) su \texttt{eth0}. La sottorete .3 \`e categoricamente bloccata verso Internet.
    \item \textbf{Accesso ai Server Pubblici da Internet}: Il traffico esterno su porte TCP 80/443 \`e indirizzato esclusivamente al Reverse Proxy in DMZ (\texttt{192.168.100.10}), il quale verifica l'URI richiesta: solo se appartiene alla whitelist \texttt{\$url1, \$url2, ...} la richiesta viene inoltrata ai server pubblici (\texttt{192.168.100.20}).
    \item \textbf{Traffico Inter-Sottorete}: Regole specifiche autorizzano l'accesso a \texttt{192.168.1.1} da .2 e .3, e l'accesso a \texttt{192.168.2.3} da .1.
\end{itemize}

\textbf{Tabella Regole del Firewall di Bordo (Default Policy: DROP):}
\begin{center}
\small
\begin{tabularx}{\textwidth}{lllllllX}
\toprule
\textbf{ID} & \textbf{Src IP} & \textbf{Dst IP} & \textbf{Proto} & \textbf{Dst Port} & \textbf{State} & \textbf{Azione \& Motivazione} \\
\midrule
1 & Any & Any & Any & Any & ESTABLISHED, RELATED & \textbf{ACCEPT}: Risposte del traffico bidirezionale \\
2 & 192.168.1.0/24 & 192.168.1.0/24 & Any & Any & Any & \textbf{ACCEPT}: Nessuna restrizione interna .1 \\
3 & 192.168.2.0/24 & 192.168.2.0/24 & Any & Any & Any & \textbf{ACCEPT}: Nessuna restrizione interna .2 \\
4 & 192.168.3.0/24 & 192.168.3.0/24 & Any & Any & Any & \textbf{ACCEPT}: Nessuna restrizione interna .3 \\
5 & 192.168.2.0/24, 192.168.3.0/24 & 192.168.1.1 & TCP/UDP & Any & NEW & \textbf{ACCEPT}: Accesso a server .1.1 da .2 e .3 \\
6 & 192.168.1.0/24 & 192.168.2.3 & TCP/UDP & Any & NEW & \textbf{ACCEPT}: Accesso a server .2.3 da .1 \\
7 & 192.168.1.0/24, 192.168.2.0/24 & Any (Internet) & Any & Any & NEW & \textbf{ACCEPT}: Uscita verso Internet per .1 e .2 (con SNAT) \\
8 & Any (Internet) & Reverse Proxy (DMZ) & TCP & 80, 443 & NEW & \textbf{ACCEPT}: Accesso al Proxy da Internet (ispezione \$url) \\
9 & Reverse Proxy & Server Pubblici & TCP & 80, 443 & NEW & \textbf{ACCEPT}: Inoltro richieste lecite post-filtro URL \\
10 & Any & Any & Any & Any & Any & \textbf{DROP (Default)}: Blocca .3 verso Internet e accessi esterni non consentiti \\
\bottomrule
\end{tabularx}
\end{center}
\end{solutionbox}

\begin{altbox}[Alternative Implementation Note: Strutture Rust per un Reverse Proxy con Filtro URL]
Se si implementa il proxy inverso di filtraggio URL in Rust:
\begin{itemize}
    \item \texttt{std::sync::Arc<HashSet<String>>}: L'insieme degli URL noti (\texttt{\$url1, \$url2}) viene memorizzato in un \texttt{HashSet} protetto da \texttt{Arc}, consentendo a centinaia di thread worker Tokio di verificare la validit\`a della richiesta HTTP in tempo medio $O(1)$ senza alcuna contesa di lock.
    \item \texttt{tokio::sync::RwLock}: Qualora la lista delle URL lecite debba essere aggiornata a caldo (senza riavviare il servizio), si utilizza un \texttt{RwLock} asincrono per permettere la sostituzione atomica dell'insieme.
\end{itemize}
\end{altbox}


\end{document}
